\chapter{International Relations}
\section{Groupings}
\section{Major Initiatives and Agreements}
\subsection{Pax-Silica Initiative (US-led)}
\begin{itemize}
    \item \textbf{Led by:} US Department of State
    \item \textbf{Objective:} Secure a resilient, trusted supply chain for \textbf{ACES} (\textbf{A}I, \textbf{C}ritical minerals, \textbf{E}lectronics \& \textbf{S}emiconductors).
    \item \textbf{India's Involvement:} India signed in during the India AI Impact Summit.
    \item \textbf{Latest Member:} Sweden
    \item \textbf{Member Nations:} Australia, Japan, ROK (Republic of Korea), UK, Singapore, Qatar, UAE, Greece, Israel, India, and Sweden.
    \item \textbf{Total Members:} 11
\end{itemize}
\section{Bilateral Relations}
\section{India and the World}
\subsection{India's West Asia Policy: Framework, Interventions, and Strategic Balancing}

\begin{itemize}
    \item \textbf{Core Pillars of Strategic Engagement:}
    \begin{enumerate}
        \item \textbf{Energy Security:} Heavy reliance on regional crude oil and gas imports.
        \item \textbf{Human Security \& Economy:} Protection of the massive Indian diaspora, safeguarding remittance inflows, and attracting sovereign wealth investments.
        \item \textbf{Geopolitics \& Trade:} Critical for maritime shipping lanes, regional connectivity nodes, and defense/counter-terrorism partnerships.
    \end{enumerate}

    \item \textbf{The "Rubik's Cube" Diplomacy (Tri-Directional Balancing):}
    \begin{enumerate}
        \item \textbf{Israel Pillar:} Anchored in core national security, defense procurement, advanced technology transfers, intelligence sharing, and counter-terrorism cooperation.
        \item \textbf{Gulf States Pillar (Saudi Arabia, UAE):} Vital economic, investment, energy, and food security partners. Requires handling Gulf countries as distinct actors rather than a single homogeneous bloc.
        \item \textbf{Iran Pillar:} Critical for continental connectivity, developing the strategic \textbf{Chabahar Port}, and maintaining alternative transit access to Afghanistan and Central Asia.
    \end{enumerate}

    \item \textbf{Emerging Regional Vulnerabilities \& Security Interconnections:}
    \begin{enumerate}
        \item \textbf{Cascading Geopolitical Crises:} Interconnected flashpoints including the Gaza conflict, US-Iran tensions, and multi-front militia activities (e.g., in Lebanon).
        \item \textbf{Maritime Security Risks:} Direct threats to trade and shipping routes in the Red Sea and Arabian Sea, causing cascading global impacts on supply lines.
    \end{enumerate}

    \item \textbf{Direct Macroeconomic and Social Impact on India:}
    \begin{enumerate}
        \item \textbf{Fiscal Strains:} Regional volatility directly triggers spikes in imported oil and fertilizer prices, depleting foreign exchange reserves and aggravating exchange-rate volatility.
        \item \textbf{Expatriate Welfare:} Instability creates physical and economic security anxieties among the Indian expatriate worker communities.
        \item \textbf{Commercial Disruptions:} Trade bottlenecks and maritime commerce delays directly affect India's export-import schedules.
    \end{enumerate}

    \item \textbf{Strategic Autonomy \& Foreign Policy Way Forward:}
    \begin{enumerate}
        \item \textbf{De-hyphenated Engagement:} Rejection of exclusive alignment with either the Israeli or Arab blocs; absolute prioritization of multi-aligned strategic autonomy.
        \item \textbf{Dialogue Maintenance:} Preserving multi-stakeholder diplomatic access while systematically avoiding strategic overcommitment to any single regional sub-rivalry.
        \item \textbf{"Access over Visibility" Doctrine:} Emphasizing high-yield backchannel diplomacy, quiet engagement, and continuous communication lines over public positioning or highly visible, polarising declarations during active conflicts.
    \end{enumerate}
\end{itemize}
\section{UN and It's agreeements}
\section{Miscellaneous}
\subsection{Medical and Wellness Tourism in India}

\begin{itemize}
    \item \textbf{ Global Market Trends:}
    \begin{enumerate}
        \item The global MVT market, valued at USD 115.6 billion in 2022, is projected to expand to USD 286.1 billion by 2030 at a CAGR of 10.8\%.
        \item India's domestic medical tourism footprint was valued at USD 8.7 billion in 2025 and is projected to scale up to USD 16.2 billion by 2030.
        \item According to the Medical Tourism Index (MTI) 2020--21, India ranks \textbf{10th} globally in medical tourism, \textbf{12th} in wellness tourism markets, and \textbf{5th} in the Asia-Pacific region for wellness.
    \end{enumerate}

    \item \textbf{The Core Pillars of India's MVT Framework: Heal In India}
    \begin{enumerate}
        \item \textbf{Medical Tourism (Curative Care):} Focuses on complex, high-end surgical procedures, organ transplants, advanced diagnostics, and clinical interventions.
        \item \textbf{Wellness Tourism (Proactive/Preventive Care):} Centered around therapeutic, preventive health via traditional systems including Yoga, Ayurveda, Naturopathy, and other AYUSH practices.
    \end{enumerate}
    \item \textbf{India's Structural and Quality Competitive Edge:}
    \begin{enumerate}
        \item \textbf{Resource Base:} Backed by 69,364 hospitals (43,486 private; 25,778 public) and 1.2 million registered doctors, fulfilling the WHO-recommended doctor-to-population ratio. English-medium clinical training removes language barriers.
        \item \textbf{Accreditation Standards:} As of 2026, the National Accreditation Board for Hospitals and Healthcare Providers (NABH) has accredited over 1,299 hospitals across 600+ safety parameters. NABH is affiliated with the International Society for Quality in Healthcare (ISQua).
        \item \textbf{Cost-Effective Medical Treatment:} India’s medical tourism appeal is significantly strengthened by its cost competitiveness, offering high-quality treatment at a substantially lower cost compared to many developed countries while maintaining comparable clinical standards.

    \item \textbf{Strengthening AYUSH-led Medical Value Travel:}
    \begin{itemize}
        \item \textbf{Unique Wellness Advantage:} India leverages its traditional systems of medicine—Ayurveda, Yoga, Naturopathy, Unani, Siddha, and Homeopathy (AYUSH)—promoting them as integral components of holistic health and preventive care.
        \item \textbf{AYUSH Visa Facilitation:} To streamline access, a dedicated AYUSH Visa category was introduced on July 27, 2023, allowing foreign nationals and their attendants to travel specifically for recognized traditional treatments.
        \item \textbf{Quality Standards:} To enhance credibility and standardization, the Bureau of Indian Standards (BIS) has officially adopted the international standard ISO 22525 for medical wellness tourism services.
        \item \textbf{Insurance Coverage:} Under IRDAI regulations, insurance coverage has expanded to AYUSH therapies, with around 27 insurance companies currently offering more than 140 policy products.
        \item \textbf{Global Engagement \& Capacity Building:} MVT remains a central theme across major forums, including the "Global Synergy in AYUSH" Summits in Mumbai (2024) and Chennai (May 2025). Workforce development is supported by the AYUSH sub-council under the Health Sector Skill Council.
        \item \textbf{Global Outreach:} Outreach continues to expand through platforms like the WHO Global Traditional Medicine Summit, the Know India Programme, and designated showcases during the Maha Kumbh Mela 2025.
    \end{itemize}
    \end{enumerate}

    \item \textbf{Budgetary and Strategic Initiatives (Union Budget 2026--27):}
    \begin{enumerate}
        \item \textbf{Regional Medical Hubs:} Plan to set up \textbf{five} Regional Medical Hubs via Public-Private Partnership (PPP) to establish integrated complexes housing healthcare, research, AYUSH facilities, and MVT Facilitation Centres under one roof.
        \item \textbf{Traditional Medicine Infrastructure:} Announcement of \textbf{three} new All India Institutes of Ayurveda. Simultaneously, the WHO Global Traditional Medicine Centre in Jamnagar is being upgraded to enhance global clinical research.
    \end{enumerate}

    \item \textbf{Policy, Governance \& Digital Facilitation:}
    \begin{enumerate}
        \item \textbf{Dedicated Visa Regime:} Introduction of specific \textbf{e-AYUSH Visa} and \textbf{e-AYUSH Attendant Visa} categories (launched July 27, 2023) alongside conventional e-Medical visas extended to 172 countries.
        \item \textbf{Standardization \& Insurance Integration:} Bureau of Indian Standards (BIS) adopted the international ISO 22525 standard for wellness tourism. IRDAI framework has successfully onboarded $\approx$27 insurance firms offering over 140 policy products covering AYUSH treatments.
        \item \textbf{Governance Bodies:} The National Medical \& Wellness Tourism Promotion Board (NMWTB) acts as an apex inter-ministerial coordinator under the Ministry of Tourism .
        \item \textbf{Capacity Building:} A pilot project under the 2026--27 Budget targets upskilling 10,000 guides at 20 iconic locations with a 12-week course. Paramedical staffs are being formally trained in cross-cultural sensibilities and foreign languages.
    \end{enumerate}
\end{itemize}